% !TeX root = ../thesisv2.tex

\chapter{绪论}
\label{cha:v2-intro}

\section{研究背景及意义}
随着高分辨率卫星、航空摄影平台和对地观测传感器持续演进，遥感影像的空间分辨率、获取频率和覆盖范围均得到显著提升。遥感影像已经从单纯的数据档案转变为城市运行监测、交通设施维护、土地利用分析、灾后快速评估和数字地图更新中的核心信息来源\cite{mnih2013,deepglobe2018}。在这一背景下，如何从海量遥感图像中快速、稳定、自动地提取道路网络，已经成为遥感图像理解和计算机视觉交叉领域的重要问题。

道路目标在遥感图像中具有明显区别于一般区域目标的结构属性。第一，道路通常呈现细长、连续、方向性明显和分支丰富的几何形态，局部宽度较小而整体跨度较大，模型需要同时兼顾局部边界和全局拓扑。第二，道路表面材质和周围地物关系复杂，建筑阴影、裸地边缘、停车场纹理和树木遮挡常导致道路与非道路区域出现显著混淆。第三，道路提取的最终使用对象通常并不是独立 patch，而是完整道路网络，因此对整图层面的连通性、稳定性和空间一致性提出了更高要求。

从应用角度看，自动道路提取能够直接服务于多种下游场景。在交通基础设施管理中，道路分布信息是地图维护、路网更新和交通状态建模的重要基础；在应急响应中，道路通行区域的快速识别可为救援调度、路径规划和风险评估提供依据；在自动驾驶先验地图建设中，高质量道路掩膜可以为高精地图生成与更新提供数据支持；在国土空间治理和智慧城市建设中，道路作为人工地表的重要组成部分，也是空间格局分析的重要基础数据。与人工判读和人工矢量化相比，基于深度学习的自动方法能够显著降低标注和更新成本，提高作业效率。

从研究角度看，遥感道路提取兼具细粒度分割与结构化建模两类特征。一方面，道路边界和细小支路对浅层纹理与空间细节高度敏感；另一方面，道路连通性和大尺度走向又依赖高层语义与上下文表达。因此，遥感道路提取不仅可以作为语义分割模型的应用场景，还能够作为检验编码器--解码器结构、跳跃连接利用方式、注意力重标定能力和多尺度上下文建模能力的重要任务平台。围绕该任务开展研究，既具有现实应用价值，也具有方法设计层面的理论意义。

\begin{table}[htbp]
\centering
\caption{遥感道路提取的主要任务特征与对应设计要求}
\label{tab:v2-task-challenges}
\begin{tabular}{p{3.2cm}p{5.2cm}p{5.2cm}}
\toprule
任务特征 & 典型表现 & 模型设计要求 \\
\midrule
细长结构明显 & 道路宽度小、跨度大、边界细 & 保持浅层边缘信息与局部纹理分辨力 \\
拓扑连续性要求高 & 主干路与支路相连，存在交叉和分叉 & 提升长程上下文建模与整图结构一致性 \\
背景干扰复杂 & 阴影、植被、裸地和停车场纹理与道路混淆 & 增强道路相关特征，抑制背景响应 \\
尺度变化显著 & 主干道路、支路和窄路共存 & 构建多尺度特征表达与融合机制 \\
使用对象为整图结果 & 需要完整道路网络而非局部 patch 结论 & 采用整图级推理和统一整图评估协议 \\
\bottomrule
\end{tabular}
\end{table}

基于上述背景与任务特征，本文将研究重点放在两个层面。其一，围绕统一 baseline 构建规范、可复核、可持续扩展的实验工程，使模型比较建立在一致的数据与评估条件之上。其二，围绕 U-Net 主干展开结构设计，一条路线聚焦跳跃连接特征增强，另一条路线聚焦瓶颈层上下文增强，并通过统一整图验证与整图测试形成数据导向结论。

\section{国内外研究现状}
\subsection{传统道路提取方法}
早期遥感道路提取方法主要依赖人工设计特征和规则驱动策略，例如边缘检测、形态学滤波、活动轮廓模型、线结构跟踪和图搜索等。这类方法通常通过灰度变化、梯度信息、线性结构先验或形态学连通规则识别道路候选区域，再通过后处理完成细化与连通。传统方法在纹理简单、背景单一或道路形态规则的场景中具有一定可解释性，但普遍依赖阈值经验和人工先验，难以适应复杂背景和大尺度场景变化。

随着遥感影像分辨率不断提升，建筑阴影、车辆遮挡、裸土边缘和植被覆盖等复杂因素显著增加，传统方法的稳定性受到较大限制。尤其是在道路宽度变化明显、交叉口结构复杂和局部遮挡严重的场景中，仅依赖局部边缘与形态规则往往难以维持道路网络的完整性。因此，研究重心逐步转向具备自动特征学习能力的深度学习方法。

\subsection{基于卷积神经网络的语义分割方法}
卷积神经网络的发展使道路提取问题逐步从人工特征识别转向端到端语义分割。U-Net 通过对称式编码器--解码器结构和跳跃连接机制，在保持高层语义表达的同时充分利用浅层细节信息，成为图像分割领域极具代表性的网络结构\cite{unet2015}。在遥感道路提取场景中，U-Net 的结构优势尤其明显：编码器能够逐步聚合全局语义，解码器结合浅层特征恢复边界和窄路细节，适合处理道路这一兼具局部细节与整体连通性的目标类型。

在 U-Net 之外，DeepLab 系列模型通过空洞卷积扩大感受野，在不显著降低分辨率的条件下增强上下文建模能力。DeepLabv3+ 将空洞卷积、编码器--解码器结构和多尺度上下文聚合结合起来，在复杂场景语义分割中表现稳定\cite{deeplabv3plus2018}。这一路线说明，上下文信息对目标与背景的区分具有显著价值，尤其适合需要长程结构表达的场景。

对于遥感道路提取任务而言，卷积神经网络方法的关键问题通常体现在三个方面。第一，浅层特征虽然包含边缘和纹理细节，但也会引入大量背景噪声；第二，高层特征具有更强语义表达能力，但在下采样过程中可能损失细小支路与边缘信息；第三，道路结构本身具有显著尺度差异与方向性，单一尺度特征往往不足以同时兼顾主干道路和窄路分支。因此，围绕跳跃连接利用方式、多尺度上下文表达和结构重标定能力的改进成为后续研究的重要方向。

\subsection{注意力机制与结构重构方法}
在 U-Net 主干基础上，研究者进一步围绕注意力机制、特征重标定和结构重构展开扩展。Attention U-Net 在跳跃连接进入解码器之前引入门控注意力，使解码器能够更有针对性地选择浅层空间区域\cite{attentionunet2018}。UNet++ 则通过嵌套式跳跃连接增加不同层之间的特征融合路径，从而缩小编码器特征与解码器特征之间的语义差距\cite{unetpp2018}。这类方法说明，浅层特征并非越多越好，关键在于以合适方式进行筛选、增强和融合。

另一类改进路线聚焦于更强的上下文表达或更大规模的网络结构。DDU-Net 通过双解码器结构和专门面向道路提取设计的上下文模块，强调高分辨率遥感场景中的多尺度道路表达\cite{ddu2022}。与此同时，空洞卷积\cite{dilated2016}和 Transformer 结构\cite{swin2021}也被引入遥感图像理解任务中，用于增强感受野、长程依赖和全局交互能力。这些方法共同说明，遥感道路提取的核心并不仅仅是逐像素分类，而是如何在高层语义、浅层细节和全局连通性之间取得合理平衡。

\subsection{现有研究对本文的启示}
结合现有研究可以得到三点直接启示。第一，U-Net 及其变体仍是遥感道路提取中最适合进行结构分析和工程复现的主干框架，原因在于其拓扑清晰、可控性强、适于系统对比。第二，跳跃连接特征增强与瓶颈层上下文增强分别对应浅层细节利用和高层上下文表达两个关键方向，均具有明确研究价值。第三，模型性能结论不能脱离统一工程条件和统一评估协议，否则不同实验之间难以直接比较。

本文正是在上述启示基础上开展工作：一方面以统一 U-Net baseline 为锚点，保证不同模型之间的比较建立在统一条件之上；另一方面同时布局主线模型、baseline 直接改进分支和扩展探索分支，使模型分析不仅停留在单一结构，而是形成层次清晰的结果矩阵。

\section{研究目标与研究内容}
\subsection{研究目标}
本文的总体目标是面向高分辨率遥感影像道路提取任务，建立一套可复用、可编译、可评估、可扩展的论文与实验一体化工程，并在该工程基础上完成以 U-Net 为中心的结构改进研究。围绕这一总体目标，本文进一步细化为以下四项具体目标：
\begin{enumerate}
  \item 建立统一的数据处理、训练、推理和整图评估流程，保证全部模型在相同实验条件下进行比较；
  \item 以参数规模适中且性能稳定的 baseline U-Net 作为统一参照，完成主线模型与改进模型的结构设计；
  \item 围绕跳跃连接特征增强和瓶颈层上下文增强两条路线开展实验，分析不同结构改动对精度、参数量和整图表现的影响；
  \item 将全部可复核实验结果组织为清晰的论文结论体系，形成能够支撑后续扩展研究的工程基础。
\end{enumerate}

\subsection{研究内容}
围绕上述目标，本文的研究内容主要包括以下几个方面。

第一，完成实验工程重建。原有材料中存在代码缺失、训练入口不完整和配置体系不统一等问题。本文基于当前本地完整材料，重新组织数据集读取、模型注册、训练器、损失函数、整图推理、整图评估和可视化导出模块，形成统一工程主线。

第二，建立统一 baseline。本文以 Batch Normalization 与双线性上采样风格的 U-Net 为全文统一 baseline，完成训练、验证和测试，并以整图级 IoU 和 F1 结果作为后续所有模型分析的统一参照。

第三，完成主线模型 DLGU-Net 设计与验证。该模型围绕跳跃连接展开，将轻量级双线性注意力模块嵌入四级跳跃连接中，以增强道路相关特征、抑制复杂背景响应，并在统一协议下完成结构对照与定量验证。

第四，开展 baseline 直接改进研究。本文在 baseline 的瓶颈层引入并行空洞卷积分支，构建空洞卷积增强基线，考察多尺度上下文增强对道路提取性能的影响，并将其作为本文围绕 baseline 展开的直接改进路线。

第五，构建扩展探索分支。为体现结构探索深度，本文进一步实现并训练 Attention U-Net、UNet++、DDU-Net、ResNet34 U-Net、Vanilla U-Net 与 Residual Vanilla U-Net 等多组模型，形成完整实验矩阵，用于丰富模型性能与结构特征的比较视角。

\subsection{技术路线}
本文的技术路线可以概括为“工程重建--统一训练--整图评估--结果重组”四个连续环节。

\begin{enumerate}
  \item 在工程准备阶段，完成数据目录重组、离线 patch 数据集生成、模型工厂设计和训练器实现，建立统一实验平台。
  \item 在模型设计阶段，围绕统一 baseline 分别构建 DLGU-Net、空洞卷积增强 baseline 以及多组补充对照模型，形成主线模型、直接改进分支和扩展探索分支。
  \item 在实验执行阶段，采用统一输入尺度、统一损失函数、统一优化器和统一整图评估协议，对全部模型完成训练、验证和测试。
  \item 在结果组织阶段，对整图验证与整图测试结果进行系统汇总，结合参数量、训练稳定性、定性结果和结构差异，形成论文主线与层次清晰的比较结论。
\end{enumerate}

\section{论文结构安排}
围绕上述研究目标和技术路线，本文将全文组织为五章。

第一章为绪论，主要说明遥感道路提取的研究背景、应用意义、国内外研究现状、本文研究目标、研究内容和技术路线，并明确全文结构安排。

第二章为理论基础、数据集与实验设置，集中介绍卷积神经网络与 U-Net 系列方法的理论基础、注意力与空洞卷积等关键机制、数据集组织方式、统一训练配置、工程优化过程以及整图评估协议，为后续模型分析提供共用基础。

第三章为基于 DLGU-Net 的遥感道路提取方法，重点阐述主线模型的研究动机、整体结构、关键模块、跳跃连接集成方式、实验结果和结构对照分析，并在统一整图协议下给出 DLGU-Net 的量化结论。

第四章为基于空洞卷积增强 baseline 的遥感道路提取方法，重点分析瓶颈层上下文增强设计、baseline 直接改进结果、效能对比实验，以及扩展探索分支中的多组模型结果，从而形成更完整的模型选择依据。

第五章为结论与展望，对全文工作进行系统总结，归纳工程层面、方法层面和实验层面的主要结论，并给出后续可继续深入的研究方向。

\section{本章小结}
本章从遥感道路提取的应用背景、任务特征和研究价值出发，说明了本文围绕 U-Net 主干开展结构设计和实验重组的必要性。在此基础上，系统梳理了传统方法、卷积神经网络方法以及注意力与结构重构方法的发展脉络，并进一步明确了本文的研究目标、研究内容、技术路线和章节安排。上述内容为后续章节展开统一理论分析、模型设计和实验比较奠定了问题背景与研究框架。
